\section{Named-Competitor Comparison Scaffold}
\label{sec:evmgpu-named-comp}
%% ============================================================================

This section pins the EVM acceleration claim to specific competitors with
their published numbers. Lux columns are \texttt{[TBD-measure]} except
where this paper has already measured the value (those cells cite the
relevant internal section).

\subsection{Competitors and Why}

\begin{itemize}
  \item \textbf{geth} (\texttt{go-ethereum})~\cite{gethgithub} ---
        Ethereum mainnet reference client; the baseline every other
        execution layer is measured against.
  \item \textbf{reth} (Paradigm)~\cite{rethgithub,paradigmreth2024} ---
        Rust execution client. Paradigm's published target is 1\,Ggas/s;
        their April 2024 post reports 100--200\,Mgas/s live sync and
        a 690\,Mgas/s historical-sync run on opBNB.
  \item \textbf{erigon} (\texttt{erigontech/erigon})~\cite{erigongithub}
        --- Go execution client with staged sync; widely reported as
        4$\times$ faster than geth on historical sync~\cite{chainstackgetherigon}.
  \item \textbf{evmone} (\texttt{ipsilon/evmone})~\cite{evmonegithub} ---
        C++ interpreter; baseline for ``how fast can the interpreter
        alone go'' (already measured in this paper).
  \item \textbf{pevm} / Block-STM (Aptos, Risechain)~\cite{blockstmppopp,pevmrise}
        --- parallel-EVM baselines. Block-STM publishes up to 170k
        TPS / 32-thread; pevm reports similar magnitudes on EVM.
\end{itemize}

\subsection{Throughput Scaffold}

\begin{table}[h]
\centering
\small
\begin{tabular}{@{}lllrl@{}}
\toprule
\textbf{Client} & \textbf{Backend} & \textbf{Workload} & \textbf{Throughput} & \textbf{Source} \\
\midrule
\textbf{Lux GPU EVM kernel}     & Metal (M1 Max)    & opcode mix       & 85.9\,M ops/s       & Tab.~\ref{tab:peak} \\
\textbf{Lux C++ evmone+state}   & CPU               & full block        & 20.9\,Ggas/s        & Tab.~\ref{tab:cpp_vs_go} \\
\textbf{Lux Go EVM (baseline)}  & CPU               & full block        & 4.2--5.4\,Ggas/s    & Tab.~\ref{tab:cpp_vs_go} \\
\textbf{Lux + parallel ecrecover} & 10-core CPU     & block w/ sigs    & 7.1$\times$ over 1-core & Tab.~\ref{tab:ecrecover} \\
\midrule
geth (latest)                   & CPU               & live mainnet sync & $\sim$25--50\,Mgas/s (typ.) & \cite{gethgithub,paradigmreth2024} \\
reth (Paradigm)                 & CPU               & live mainnet sync & 100--200\,Mgas/s    & \cite{paradigmreth2024} \\
reth (Paradigm)                 & CPU               & historical opBNB sync & 690\,Mgas/s        & \cite{paradigmreth2024} \\
erigon                          & CPU               & full sync vs.\ geth & 4$\times$ geth      & \cite{chainstackgetherigon} \\
evmone (C++ interp.)            & CPU               & interp.-only      & $\sim$416\,Ggas/s   & this paper, Tab.~\ref{tab:cpp_vs_go} \\
Block-STM (Aptos)               & 32-thread CPU     & Aptos workload    & up to 170k TPS      & \cite{blockstmppopp} \\
Block-STM (Diem)                & 32-thread CPU     & Diem workload     & up to 110k TPS      & \cite{blockstmppopp} \\
\bottomrule
\end{tabular}
\caption{EVM/execution throughput comparison scaffold. Lux rows are
already measured (this paper); the comparison is against the cited
upstream-published numbers. The units differ deliberately (TPS vs.\ Mgas/s
vs.\ opcodes/s) --- the Reproducibility spec below unifies them.}
\label{tab:evmgpu-comparison}
\end{table}

\subsection{Crypto Bottleneck Comparison}

\begin{table}[h]
\centering
\small
\begin{tabular}{@{}lrrl@{}}
\toprule
\textbf{Operation} & \textbf{Lux (this paper)} & \textbf{Baseline} & \textbf{Source} \\
\midrule
\texttt{ecrecover} (single)       & \texttt{[TBD-measure]}              & $\sim$60--100\,\textmu s (\texttt{secp256k1})  & \cite{secp256k1lib} \\
\texttt{ecrecover} parallel       & 7.1$\times$ on 10 cores                                  & linear scaling (lib-limited) & Tab.~\ref{tab:ecrecover} \\
\texttt{ecrecover} GPU (Metal)    & 20--24$\times$ vs.\ CPU (Lux memory)                     & none published & \texttt{[TBD-baseline]} \\
BLS12-381 aggregate verify        & \texttt{[TBD-measure]}                                   & blst 1.0--1.5\,ms (n$\sim$128) & \cite{supranationalblstevm} \\
\bottomrule
\end{tabular}
\caption{Crypto-operation comparison. \texttt{ecrecover} dominates 87\%
of Go EVM block-processing time (Sec.~5, Tab.~\ref{tab:profile}), so this
is the highest-leverage row to populate.}
\label{tab:evmgpu-crypto}
\end{table}

\subsection{Reproducibility Specification}
\label{sec:evmgpu-repro}

To turn the unit-normalised entries in Table~\ref{tab:evmgpu-comparison}
into apples-to-apples cells, all five backends must be re-run on the
same machine against the same trace.

\begin{itemize}
  \item \textbf{Hardware.} (1) Apple M1 Max 64\,GiB (matches this paper's
        existing measurements); (2) AWS \texttt{c7i.16xlarge} (Intel
        Sapphire Rapids, 64 vCPU); (3) AWS \texttt{p4d.24xlarge} (8$\times$A100)
        for the GPU path.
  \item \textbf{Workload.} Replay mainnet blocks 17{,}000{,}000 --
        17{,}010{,}000 (10k blocks, $\sim$30\,M gas avg, $\sim$160\,M
        tx-sigs total). This is the same range Paradigm used in
        \cite{paradigmreth2024}.
  \item \textbf{Builds.} \texttt{ethereum/go-ethereum@latest};
        \texttt{paradigmxyz/reth@latest}; \texttt{erigontech/erigon@latest};
        \texttt{ipsilon/evmone@latest}; \texttt{risechain/pevm@latest}.
  \item \textbf{Metric normalisation.} Report Mgas/s for every backend;
        TPS only as a secondary metric (since gas-per-tx varies). Walltime
        of full replay; CPU\%/thread; peak RSS; disk read GB.
  \item \textbf{Commands.}
  \begin{verbatim}
    # reth (sync from snapshot to block 17_010_000)
    reth import --chain mainnet ./mainnet-17M.rlp
    # geth (import RLP)
    geth import ./mainnet-17M.rlp
    # erigon (stage sync to specific block)
    erigon --chain mainnet --datadir ./erigon-data
    # Lux (this paper)
    bench-evmgpu --replay ./mainnet-17M.rlp --backend gpu
  \end{verbatim}
  \item \textbf{Decision rule.} Lux GPU EVM wins only if Mgas/s on the
        Apple M1 Max GPU path \emph{exceeds reth on a comparable-class
        CPU machine}, measured on the same RLP range. The current
        4.85$\times$ headline is opcodes/s, not gas/s --- the
        gas/s number on the same trace is the falsifiable claim.
\end{itemize}

%% ============================================================================
